\documentclass[11pt]{article}
\usepackage[a4paper,margin=1in]{geometry}
\usepackage{fontspec}
\usepackage[boldfont=false]{xeCJK}
\setCJKmainfont{FandolSong-Regular.otf}
\usepackage{amsmath}
\usepackage{amssymb}
\usepackage{array}
\usepackage{booktabs}
\usepackage{graphicx}
\usepackage{adjustbox}
\usepackage{enumitem}
\setlist[itemize]{topsep=2pt,partopsep=0pt,parsep=0pt,itemsep=2pt,leftmargin=1.4em}
\setlist[enumerate]{topsep=2pt,partopsep=0pt,parsep=0pt,itemsep=2pt,leftmargin=1.6em}
\usepackage{parskip}
\usepackage{fvextra}
\fvset{breaklines=true,breakanywhere=true,fontsize=\small}
\setlength{\parindent}{0pt}

\begin{document}

\begin{center}
  {\LARGE\bfseries Basic economic ideas and resource allocation}\\[0.35em]
  {\large A Level Economics}
\end{center}
\vspace{0.6em}

\section*{The fundamental economic problem}

Economics starts with one big idea: \textbf{scarcity} 稀缺性. We never have enough to give everyone everything they want.

\begin{itemize}
\item \textbf{resources} 资源 are the things we use to make goods and services. They are \textbf{finite} 有限的 — there is only a limited amount, so they are \textbf{scarce} 稀缺.

\item \textbf{goods} 物品 (things you can touch, like food) and \textbf{services} 服务 (work done for you, like a haircut) are what we make from resources.

\item human \textbf{wants} 欲望 are unlimited. People always want more.

\end{itemize}

Because wants are bigger than resources, this is called the \textbf{economic problem} 经济问题. Every society must answer three questions:

\begin{enumerate}
\item \textbf{What} to produce?

\item \textbf{How} to produce it?

\item \textbf{For whom} to produce it?

\end{enumerate}

\subsection*{Choice and opportunity cost}

Because resources are scarce, you must choose. Every choice means you give up something else.

\textbf{opportunity cost} 机会成本 — the next best \textbf{alternative} 替代选择 you give up when you make a decision.

Opportunity cost is faced by everyone:

\begin{itemize}
\item \textbf{consumers} 消费者 (people who buy and use goods): buy a phone, and you give up the holiday that money could have paid for.

\item \textbf{producers} 生产者 and \textbf{firms} 企业 (the businesses that make goods): use a factory to make cars, and you give up the trucks it could have made.

\item \textbf{governments} 政府: spend tax money on hospitals, and you give up the schools that money could have built.

\end{itemize}

\section*{Thinking like an economist}

\subsection*{Economics as a social science}

Economics is a \textbf{social science} 社会科学. It studies how people behave when they make choices. Like other sciences, it builds \textbf{models} 模型 (simple pictures of the real world) and tests them with data.

The real world is complex, so economists change one thing at a time and keep everything else fixed. This is the idea of \textbf{ceteris paribus} 其他条件不变 — a Latin phrase meaning "all other things stay the same". For example: "if the price of coffee rises, people buy less coffee, \emph{ceteris paribus}".

\subsection*{Positive and normative statements}

\begin{itemize}
\item a \textbf{positive statement} 实证表述 is about fact. It can be tested as true or false. Example: "a higher price of petrol reduces the amount people buy."

\item a \textbf{normative statement} 规范表述 is about opinion. It says what \emph{should} happen, and it cannot be tested as true or false. Example: "the government should make petrol cheaper."

\end{itemize}

Normative statements rest on a \textbf{value judgement} 价值判断 — a personal view about what is good or fair. Two people can agree on the facts and still disagree about what \emph{should} be done.

\subsection*{Microeconomics and macroeconomics}

\begin{itemize}
\item \textbf{microeconomics} 微观经济学 studies single parts of the economy: one market, one firm, one consumer.

\item \textbf{macroeconomics} 宏观经济学 studies the whole economy together: total \textbf{output} 产出, jobs, and rising prices.

\end{itemize}

\section*{Factors of production}

The \textbf{factors of production} 生产要素 are the four kinds of resource used to make goods and services. Each one earns a \textbf{reward} 报酬.

\begin{center}
\adjustbox{max width=\linewidth}{%
\begin{tabular}{lll}
\toprule
\textbf{Factor} & \textbf{What it is} & \textbf{Its reward} \\
\midrule
\textbf{land} 土地 & natural resources (fields, water, oil, minerals) & \textbf{rent} 地租 \\
\textbf{labour} 劳动 & human work, both physical and mental & \textbf{wages} 工资 \\
\textbf{capital} 资本 & man-made things used to produce, like machines, tools and factories & \textbf{interest} 利息 \\
\textbf{enterprise} 企业家才能 & the skill of bringing the other three together and taking the risk & \textbf{profit} 利润 \\
\bottomrule
\end{tabular}}
\end{center}

Land includes some \textbf{non-renewable resources} 不可再生资源 (like oil and coal) that can run out. Enterprise is supplied by the \textbf{entrepreneur} 企业家, the person who has the business idea and risks their own money.

\subsection*{Specialisation and division of labour}

\textbf{specialisation} 专业化 is when a worker, firm or country makes only the things it is best at, then trades for the rest.

Inside a workplace, specialisation leads to the \textbf{division of labour} 劳动分工: a big job is split into small tasks, and each \textbf{worker} 工人 does one task again and again.

\begin{center}
\adjustbox{max width=\linewidth}{%
\begin{tabular}{ll}
\toprule
\textbf{Advantages} & \textbf{Disadvantages} \\
\midrule
workers get faster and better at one task & doing one task is boring, so workers may quit \\
less time wasted moving between jobs & if one worker is absent, the whole line can stop \\
output per worker rises, so costs fall & workers learn only one narrow skill \\
\bottomrule
\end{tabular}}
\end{center}

\subsection*{Mobility of factors}

\textbf{mobility} 流动性 means how easily a factor can move to a new use.

\begin{itemize}
\item \textbf{occupational mobility} 职业流动性 — how easily a factor moves from one \emph{job or use} to another. A teacher who retrains as a nurse is occupationally mobile.

\item \textbf{geographical mobility} 地理流动性 — how easily a factor moves from one \emph{place} to another. Land cannot move at all, so it has no geographical mobility.

\end{itemize}

Labour is often not very mobile. Family ties, the cost of housing, and missing skills all make it hard for workers to move.

\section*{Economic systems}

An economic system is the way a country decides how to \textbf{allocate} 配置 (share out) its scarce resources. There are three types.

\begin{itemize}
\item a \textbf{market economy} 市场经济 — firms and people own resources, and prices decide what is made. The government does little.

\item a \textbf{planned economy} 计划经济 (also called a \textbf{command economy} 指令经济) — the government owns resources and decides what is made, how, and for whom.

\item a \textbf{mixed economy} 混合经济 — a blend of the two. The \textbf{private sector} 私营部门 (firms and households) and the \textbf{public sector} 公共部门 (the government) both make decisions. Almost every real country is mixed.

\end{itemize}

\subsection*{The price mechanism}

In a market economy, resources are guided by the \textbf{price mechanism} 价格机制. Prices do three jobs:

\begin{itemize}
\item \textbf{signalling} 信号传递 — a price tells buyers and sellers where resources are wanted. A rising price says "this is wanted — make more".

\item \textbf{incentive} 激励 — a higher price rewards producers, so they want to supply more of that good.

\item \textbf{rationing} 配给 — when a good is scarce, a higher price shares it out to those who are willing and able to pay.

\end{itemize}

\subsection*{Comparing the systems}

\begin{center}
\adjustbox{max width=\linewidth}{%
\begin{tabular}{lll}
\toprule
\textbf{System} & \textbf{Strengths} & \textbf{Weaknesses} \\
\midrule
market & choice, low prices, strong incentive to work and invent & large gap between rich and poor; some goods not supplied \\
planned & smaller gap between rich and poor; key goods provided for all & weak incentive; shortages and surpluses; slow to change \\
mixed & tries to keep the good points of both & the right balance is hard to find \\
\bottomrule
\end{tabular}}
\end{center}

\section*{Production possibility curves}

A \textbf{production possibility curve} 生产可能性曲线 (PPC) is a line showing the largest combinations of two goods an economy can make when all its resources are fully and well used.

Imagine an economy that makes only food and machines. The PPC is drawn with food on one axis and machines on the other.

\begin{itemize}
\item a point \textbf{on} the curve: all resources are used with full \textbf{efficiency} 效率.

\item a point \textbf{inside} the curve: resources are wasted (for example, \textbf{unemployment} 失业 of workers). The economy could do better.

\item a point \textbf{outside} the curve: impossible today, because there are not enough resources. This shows scarcity.

\end{itemize}

Moving along the curve shows opportunity cost: to make more food, you must move resources away from machines, so you make fewer machines. This is the \textbf{trade-off} 取舍 between the two goods.

The \emph{shape} of the curve shows how opportunity cost behaves:

\begin{itemize}
\item a straight-line PPC means \textbf{constant} opportunity cost — each extra unit of food costs the same number of machines.

\item a curve that bows outwards means \textbf{increasing} opportunity cost — each extra unit of food costs more and more machines, because resources are not equally good at making both goods.

\end{itemize}

\subsection*{Shifts and growth}

When the PPC moves outward (to the right), the economy can make more of both goods. This is \textbf{economic growth} 经济增长. It happens when there are more resources, or better ones (new \textbf{technology} 技术, a bigger or more skilled workforce). War or disaster can move the PPC inward — this is \textbf{economic decline} 经济衰退.

A country must also choose between two kinds of good now:

\begin{itemize}
\item \textbf{capital goods} 资本品 — goods used to make other goods (machines, factories).

\item \textbf{consumer goods} 消费品 — goods people use directly (food, clothes).

\end{itemize}

Making more capital goods today means fewer consumer goods today. But more capital goods help the economy grow faster, so the PPC shifts out by more in the future.

\section*{Types of goods and services}

\subsection*{Free goods and economic goods}

\begin{itemize}
\item a \textbf{free good} 免费物品 has no opportunity cost. Its supply is unlimited, so it has no price (for example, air or sunlight).

\item an \textbf{economic good} 经济物品 (also called a \textbf{private good} 私人物品) is scarce. Making it has an opportunity cost, so it has a price.

\end{itemize}

A private good has two features. It has \textbf{rivalry} 竞争性 (if I eat the apple, you cannot) and \textbf{excludability} 排他性 (the seller can stop you using it unless you pay).

\subsection*{Public goods}

A \textbf{public good} 公共物品 is the opposite of a private good. It has two features:

\begin{itemize}
\item \textbf{non-rivalry} 非竞争性 — one person using it does not reduce the amount left for others.

\item \textbf{non-excludability} 非排他性 — you cannot stop people who do not pay from using it.

\end{itemize}

Because of non-excludability, public goods suffer from the \textbf{free-rider problem} 搭便车问题: people enjoy the good without paying for it. So private firms cannot make a profit and will not supply it. The government must provide public goods instead — for example, street lighting and \textbf{national defence} 国防.

\subsection*{Merit and demerit goods}

\begin{itemize}
\item a \textbf{merit good} 优值品 is better for you (and for society) than you realise, so people buy too little of it. Examples: education and health care.

\item a \textbf{demerit good} 劣值品 is worse for you than you realise, so people buy too much of it. Examples: cigarettes and alcohol.

\end{itemize}

People misjudge these goods because of \textbf{information failure} 信息失灵 — they do not have, or do not use, the full facts about the benefit or harm.

Deciding what counts as a merit or demerit good is partly a value judgement, so people may disagree about which goods belong on each list.


\end{document}
